% Generated by figures/06_sentiment_stats.R -- do not edit by hand.
% Requires \usepackage{booktabs} and \usepackage{threeparttable}.
\begin{table}[htbp]
\centering
\begin{threeparttable}
\caption{Performance of the hybrid sentiment classifier under three train-test approaches}
\label{tab:sentiment-stats}
\small
\begin{tabular}{l rrrrr @{\hspace{1.5em}} rrrrr}
\toprule
 & \multicolumn{5}{c}{China test data} & \multicolumn{5}{c}{US test data} \\
\cmidrule(lr){2-6} \cmidrule(lr){7-11}
\multicolumn{11}{l}{\textit{A. Three classes}} \\[2pt]
 & \multicolumn{3}{c}{F1} & Acc. & & \multicolumn{3}{c}{F1} & Acc. & \\
\cmidrule(lr){2-4} \cmidrule(lr){7-9}
 & Neg. & Amb. & Pos. & (\%) & $n$ & Neg. & Amb. & Pos. & (\%) & $n$ \\
\midrule
Approach 1 & 0.90 & 0.62 & 0.75 & 75.9 & 83 & 0.44 & 0.46 & 0.41 & 44.6 & 130 \\
Approach 2 & 0.79 & 0.40 & 0.76 & 69.9 & 73 & 0.84 & 0.39 & 0.63 & 68.6 & 70 \\
Approach 3 & 0.76 & 0.63 & 0.68 & 68.4 & 95 & 0.58 & 0.62 & 0.47 & 57.9 & 107 \\
\midrule
\multicolumn{11}{l}{\textit{B. One-vs-rest AUC and positive/negative accuracy}} \\[2pt]
 & \multicolumn{3}{c}{AUC} & Acc. & & \multicolumn{3}{c}{AUC} & Acc. & \\
\cmidrule(lr){2-4} \cmidrule(lr){7-9}
 & Neg. & Amb. & Pos. & (\%) & $n$ & Neg. & Amb. & Pos. & (\%) & $n$ \\
\midrule
Approach 1 & 0.96 & 0.56 & 0.91 & 91.4 & 58 & 0.64 & 0.50 & 0.75 & 80.0 & 55 \\
Approach 2 & 0.88 & 0.49 & 0.89 & 87.9 & 58 & 0.88 & 0.62 & 0.89 & 85.5 & 55 \\
Approach 3 & 0.96 & 0.56 & 0.87 & 92.9 & 56 & 0.74 & 0.54 & 0.80 & 78.0 & 50 \\
\bottomrule
\end{tabular}
\begin{tablenotes}[flushleft]
\footnotesize
\item \textit{Notes:} The hybrid score averages the SVM label, the random forest label and the standardized Bing Liu lexicon score.
Approach 1 trains on positive and negative documents only; approach 2 adds a sample of 75 ambiguous documents;
approach 3 trains separate positive-vs-not and negative-vs-not models and combines them.
Panel A: the continuous hybrid score is cut into three classes at quantiles matching the share of each hand-coded class in the evaluation set;
for approach 1, which has no ambiguous documents in its test split, held-out ambiguous documents are added to the evaluation set.
Panel B: AUC treats each class in turn as the outcome and the hybrid score as the predictor, with the direction of the score chosen for each class (the pROC default);
accuracy classifies hand-coded positive and negative documents by the sign of the hybrid score.
$n$ is the number of evaluation documents.
\end{tablenotes}
\end{threeparttable}
\end{table}
